(1) 理想Mach–Zehnder干涉仪的两个无损分束器用实正交矩阵 $R(\theta)=\left(\begin{smallmatrix}\cos\theta&-\sin\theta\\\sin\theta&\cos\theta\end{smallmatrix}\right)$ 表示，两角分别为 $\theta_1,\theta_2\in(0,\pi/2)$，两臂间相位差为 $\phi$，仅输入端口1有单位归一振幅。求输出端口振幅。

(2) 求两端口光强及总强度守恒；写出端口1的条纹可见度，找出可见度为1的耦合器角度关系。

(3) 求出现完全暗端口的相位及两耦合器匹配条件；在任意非对称角度下能否同时把全部光送到一个端口？

(4) 以独立单光子探测每次得到端口1或2为观测，求相位的二元Fisher信息，在两个50:50分束器极限下给出其值。
